\documentclass[UTF8, a4paper]{article}

\usepackage{geometry}
\geometry{a4paper, margin=2.5cm}

\usepackage{amsmath, amssymb, amsthm}
\usepackage{algorithm}
\usepackage{algpseudocode}
\usepackage{booktabs}
\usepackage{xcolor}
\usepackage{listings}
\lstset{
    basicstyle=\ttfamily\small,
    keywordstyle=\color{blue}\bfseries,
    commentstyle=\color{green!50!black}\itshape,
    stringstyle=\rmfamily\upshape,
    numbers=left,
    numberstyle=\tiny\color{gray},
    stepnumber=1,
    numbersep=5pt,
    backgroundcolor=\color{white},
    showstringspaces=false,
    breaklines=true,
    frame=single,
    rulecolor=\color{black},
}

\theoremstyle{definition}
\newtheorem{lemma}{Lemma}
\newtheorem{theorem}{Theorem}

\title{\bf{Variants of FFT and its applications}}
\author{Yanqiao Chen\\
        Southern University of Science and Technology\\
        Department of Computer Science and Engineering\\
        \texttt{12412115@mail.sustech.edu.cn}}
\date{\today}

\begin{document}

\maketitle

\begin{abstract}
In this assignment report, we will first recall the FFT we have learnt in class, and then introduce 
the Cooley-Tukey FFT algorithm, which is an in-place iterative implementation of the FFT.
We will also introduce the Number Theoretic Transform (NTT), which is a variant of the 
FFT that operates in a finite field. We will discuss some applications of FFT in various
 fields such as signal processing, image processing, and polynomial multiplication.
\end{abstract}

\section{Introducing Fast Fourier Transform (FFT)} 

\subsection{The Convolution} 

In many scenarios, we need to compute the convolution of two sequences. The convolution of two sequences \( f \) and \( g \) is defined as:
\[
(f * g)[n] = \sum_{m=-\infty}^{\infty} f[m] \cdot g[n - m]
\]

Trivally, computing the convolution directly has a time complexity of \( O(n^2) \). However, using the Fast Fourier Transform (FFT), 
we can compute the convolution efficiently in \( O(n \log n) \) time. For any convolution problem, we follow the steps below:

\begin{enumerate} 
    \item Encode the sequences \( f \) and \( g \) into polynomials, where the coefficients of the polynomials correspond to the elements of the sequences.
    \item Compute the FFT of both polynomials to obtain their values at the roots of unity.
    \item Perform point-wise multiplication of the FFT results to get the FFT of the convolution.
    \item Compute the inverse FFT to obtain the coefficients of the resulting polynomial, which correspond to the convolution of the original sequences.
\end{enumerate}

\subsection{FFT} 

We now recall the process of computing the FFT. Our target is to compute a vector that contains the values of a polynomial \( P(x) \) at the \( n \)-th roots of unity. The FFT can be computed using
a divide-and-conquer approach, which breaks down the problem into smaller subproblems. We notice that 
$$
w^{k + n/2}_{n} = -w^{k}_{n/2}
$$
according to the periodicity of the roots of unity. Therefore, we can express the FFT of \( P(x) \) as:
\[
P(w^k_n) = P_{\text{even}}(w^{k}_{n/2}) + w^k_n P_{\text{odd}}(w^{k}_{n/2})
\]
and 
\[
P(w^{k + n/2}_n) = P_{\text{even}}(w^{k}_{n/2}) - w^k_n P_{\text{odd}}(w^{k}_{n/2})
\]

We may do this algorithm recursively until we reach the base case where the polynomial has only one coefficient.
The time complexity of this approach is \( O(n \log n) \) due to the logarithmic depth of the recursion and the linear work done at each level.

\subsection{Recursive FFT Algorithm}

\subsubsection{Algorithm Description}

The recursive FFT follows the divide-and-conquer paradigm:

\begin{enumerate}
    \item \textbf{Base case}: If \( n = 1 \), return the single coefficient \([a_0]\).
    \item \textbf{Divide}: Split the polynomial into even and odd indexed coefficients:
    \[
    P_{\text{even}}(x) = a_0 + a_2x + a_4x^2 + \cdots + a_{n-2}x^{n/2-1}
    \]
    \[
    P_{\text{odd}}(x) = a_1 + a_3x + a_5x^2 + \cdots + a_{n-1}x^{n/2-1}
    \]
    \item \textbf{Conquer}: Recursively compute \(\text{FFT}(P_{\text{even}}, n/2)\) and \(\text{FFT}(P_{\text{odd}}, n/2)\).
    \item \textbf{Combine}: Merge results using the butterfly operation with twiddle factors.
\end{enumerate}

\subsubsection{Step-by-Step Example: Recursive FFT for \( n=8 \)}

Given coefficients \([a_0, a_1, a_2, a_3, a_4, a_5, a_6, a_7]\) and primitive root \(\omega_8 = e^{-2\pi i / 8}\).

\textbf{Level 0: Initial Call} (\( n=8 \))

Split into:
\begin{itemize}
    \item \( P_{\text{even}} = [a_0, a_2, a_4, a_6] \) (even indices)
    \item \( P_{\text{odd}} = [a_1, a_3, a_5, a_7] \) (odd indices)
\end{itemize}

\textbf{Level 1: Recursive Calls} (\( n=4 \))

For \( P_{\text{even}} = [a_0, a_2, a_4, a_6] \) with \(\omega_4 = e^{-2\pi i / 4}\):
\begin{itemize}
    \item \( P_{\text{even,even}} = [a_0, a_4] \) (indices 0, 4)
    \item \( P_{\text{even,odd}} = [a_2, a_6] \) (indices 2, 6)
\end{itemize}

For \( P_{\text{odd}} = [a_1, a_3, a_5, a_7] \) with \(\omega_4 = e^{-2\pi i / 4}\):
\begin{itemize}
    \item \( P_{\text{odd,even}} = [a_1, a_5] \) (indices 1, 5)
    \item \( P_{\text{odd,odd}} = [a_3, a_7] \) (indices 3, 7)
\end{itemize}

\textbf{Level 2: Base Cases and Combining} (\( n=2 \))

For \( P_{\text{even,even}} = [a_0, a_4] \) with \(\omega_2 = e^{-2\pi i / 2} = -1\):
\begin{itemize}
    \item Base case: \(\text{FFT}([a_0], 1) = [a_0]\), \(\text{FFT}([a_4], 1) = [a_4]\)
    \item Combine:
    \begin{align*}
    y_0 &= a_0 + \omega_2^0 \cdot a_4 = a_0 + a_4 \\
    y_1 &= a_0 - \omega_2^0 \cdot a_4 = a_0 - a_4
    \end{align*}
    \item Result: \([a_0 + a_4, a_0 - a_4]\)
\end{itemize}

Similarly for \( P_{\text{even,odd}} = [a_2, a_6] \):
\begin{itemize}
    \item Result: \([a_2 + a_6, a_2 - a_6]\)
\end{itemize}

Combine for \( P_{\text{even}} \) (\( n=4 \)) with \(\omega_4^0 = 1\), \(\omega_4^1 = -i\):
\begin{align*}
y_0 &= (a_0 + a_4) + \omega_4^0 \cdot (a_2 + a_6) = a_0 + a_4 + a_2 + a_6 \\
y_1 &= (a_0 - a_4) + \omega_4^1 \cdot (a_2 - a_6) = a_0 - a_4 - i(a_2 - a_6) \\
y_2 &= (a_0 + a_4) - \omega_4^0 \cdot (a_2 + a_6) = a_0 + a_4 - a_2 - a_6 \\
y_3 &= (a_0 - a_4) - \omega_4^1 \cdot (a_2 - a_6) = a_0 - a_4 + i(a_2 - a_6)
\end{align*}

Similarly compute for \( P_{\text{odd}} \):
\begin{align*}
y_4 &= a_1 + a_5 + a_3 + a_7 \\
y_5 &= a_1 - a_5 - i(a_3 - a_7) \\
y_6 &= a_1 + a_5 - a_3 - a_7 \\
y_7 &= a_1 - a_5 + i(a_3 - a_7)
\end{align*}

\textbf{Level 3: Final Combine} (\( n=8 \))

Combine results with \(\omega_8^k = e^{-2\pi i k / 8}\):
\begin{align*}
Y_0 &= y_0 + \omega_8^0 \cdot y_4 \\
Y_1 &= y_1 + \omega_8^1 \cdot y_5 \\
Y_2 &= y_2 + \omega_8^2 \cdot y_6 \\
Y_3 &= y_3 + \omega_8^3 \cdot y_7 \\
Y_4 &= y_0 - \omega_8^0 \cdot y_4 \\
Y_5 &= y_1 - \omega_8^1 \cdot y_5 \\
Y_6 &= y_2 - \omega_8^2 \cdot y_6 \\
Y_7 &= y_3 - \omega_8^3 \cdot y_7
\end{align*}

\subsubsection{Recursive FFT Pseudocode}

\begin{verbatim}
RecursiveFFT(a, n, omega_n):
    // a: coefficient array
    // n: length (must be power of 2)
    // omega_n: primitive n-th root of unity

    if n == 1:
        return [a[0]]

    // Split into even and odd indices
    even = [a[0], a[2], a[4], ..., a[n-2]]
    odd  = [a[1], a[3], a[5], ..., a[n-1]]

    // Recursive calls
    omega_half = omega_n * omega_n  // primitive (n/2)-th root
    y_even = RecursiveFFT(even, n/2, omega_half)
    y_odd  = RecursiveFFT(odd,  n/2, omega_half)

    // Combine using butterfly operations
    y = array of size n
    omega = 1  // omega_n^0
    for k = 0 to n/2 - 1:
        y[k]       = y_even[k] + omega * y_odd[k]
        y[k + n/2] = y_even[k] - omega * y_odd[k]
        omega = omega * omega_n  // next twiddle factor

    return y
\end{verbatim}

\subsubsection{Complexity Analysis}

\begin{itemize}
    \item \textbf{Time Complexity}: \( T(n) = 2T(n/2) + O(n) = O(n \log n) \)
    \item \textbf{Space Complexity}: \( O(n \log n) \) due to recursive call stack and temporary arrays at each level
\end{itemize}

However, the recursive implementation of FFT can lead to significant overhead due to function calls and stack usage. To optimize this,
 we can implement the FFT iteratively using an in-place algorithm.

\section{In-Place FFT Algorithm (Iterative Implementation)}

Given a polynomial $A(x) = a_0 + a_1x + a_2x^2 + \cdots + a_{n-1}x^{n-1}$ with $n = 2^k$ coefficients.

\subsection{Step 1: Bit-Reversal Permutation}

Rearrange the coefficient array in bit-reversed order. For each index $i$, compute its $k$-bit reversal $\text{rev}(i)$ and swap $a_i$ with $a_{\text{rev}(i)}$ (avoiding duplicate swaps when $i < \text{rev}(i)$).

For example, with $n=8$ (3 bits):
\begin{center}
\begin{tabular}{c|c|c}
Original Index \& Binary \& Bit-Reversed Index \\
\hline
0 \& 000 \& 0 \\
1 \& 001 \& 4 \\
2 \& 010 \& 2 \\
3 \& 011 \& 6 \\
4 \& 100 \& 1 \\
5 \& 101 \& 5 \\
6 \& 110 \& 3 \\
7 \& 111 \& 7 \\
\end{tabular}
\end{center}

The permuted array is:
$$[a_0, a_4, a_2, a_6, a_1, a_5, a_3, a_7]$$

\subsection{Step 2: Iterative Butterfly Operations}

Starting with $m=2$, double $m$ each iteration until $m=n$. For each stage, we use twiddle factors:
$$\omega_m^j = e^{-2\pi i j / m} = \cos\left(\frac{2\pi j}{m}\right) - i\sin\left(\frac{2\pi j}{m}\right)$$

\subsubsection{Stage 1: $m=2$ (groups of size 2)}

Process adjacent pairs with butterfly operations:
\begin{align*}
\ &[a_0 + a_4, \; a_0 - a_4, \; a_2 + a_6, \; a_2 - a_6, \; a_1 + a_5, \; a_1 - a_5, \; a_3 + a_7, \; a_3 - a_7] \\
\ &= [A_0^{(1)}, A_4^{(1)}, A_2^{(1)}, A_6^{(1)}, A_1^{(1)}, A_5^{(1)}, A_3^{(1)}, A_7^{(1)}]
\end{align*}

Note: $\omega_2^0 = 1$, so this stage uses only additions and subtractions.

\subsubsection{Stage 2: $m=4$ (groups of size 4)}

Twiddle factors: $\omega_4^0 = 1$, $\omega_4^1 = -i$

For the first half $[A_0^{(1)}, A_4^{(1)}, A_2^{(1)}, A_6^{(1)}]$:
\begin{align*}
A_0^{(2)} &= A_0^{(1)} + \omega_4^0 \cdot A_2^{(1)} \\
A_2^{(2)} &= A_0^{(1)} - \omega_4^0 \cdot A_2^{(1)} \\
A_4^{(2)} &= A_4^{(1)} + \omega_4^1 \cdot A_6^{(1)} \\
A_6^{(2)} &= A_4^{(1)} - \omega_4^1 \cdot A_6^{(1)}
\end{align*}

For the second half $[A_1^{(1)}, A_5^{(1)}, A_3^{(1)}, A_7^{(1)}]$:
\begin{align*}
A_1^{(2)} &= A_1^{(1)} + \omega_4^0 \cdot A_3^{(1)} \\
A_3^{(2)} &= A_1^{(1)} - \omega_4^0 \cdot A_3^{(1)} \\
A_5^{(2)} &= A_5^{(1)} + \omega_4^1 \cdot A_7^{(1)} \\
A_7^{(2)} &= A_5^{(1)} - \omega_4^1 \cdot A_7^{(1)}
\end{align*}

\subsubsection{Stage 3: $m=8$ (groups of size 8, final stage)}

Twiddle factors: $\omega_8^0 = 1$, $\omega_8^1 = \frac{\sqrt{2}}{2} - i\frac{\sqrt{2}}{2}$, $\omega_8^2 = -i$, $\omega_8^3 = -\frac{\sqrt{2}}{2} - i\frac{\sqrt{2}}{2}$

For all 8 elements:
\begin{align*}
A_0^{(3)} &= A_0^{(2)} + \omega_8^0 \cdot A_1^{(2)} \\
A_1^{(3)} &= A_0^{(2)} - \omega_8^0 \cdot A_1^{(2)} \\
A_2^{(3)} &= A_2^{(2)} + \omega_8^1 \cdot A_3^{(2)} \\
A_3^{(3)} &= A_2^{(2)} - \omega_8^1 \cdot A_3^{(2)} \\
A_4^{(3)} &= A_4^{(2)} + \omega_8^2 \cdot A_5^{(2)} \\
A_5^{(3)} &= A_4^{(2)} - \omega_8^2 \cdot A_5^{(2)} \\
A_6^{(3)} &= A_6^{(2)} + \omega_8^3 \cdot A_7^{(2)} \\
A_7^{(3)} &= A_6^{(2)} - \omega_8^3 \cdot A_7^{(2)}
\end{align*}

The final result $[A_0^{(3)}, A_1^{(3)}, A_2^{(3)}, A_3^{(3)}, A_4^{(3)}, A_5^{(3)}, A_6^{(3)}, A_7^{(3)}]$ is the DFT.

\subsection{General In-Place FFT Pseudocode}

\begin{verbatim}
FFT(a, n):
    // Step 1: Bit-reversal permutation
    for i = 0 to n-1:
        j = bit_reverse(i, log2(n))
        if i < j:
            swap(a[i], a[j])

    // Step 2: Iterative butterfly operations
    for m = 2; m <= n; m *= 2:        // m = current group size
        omega_m = exp(-2*pi*i / m)    // principal twiddle factor
        for k = 0; k < n; k += m:     // iterate over groups
            omega = 1                  // omega_m^0
            for j = 0; j < m/2; j++:  // butterfly pairs within group
                u = a[k + j]           // upper butterfly branch
                v = omega * a[k + j + m/2]  // lower butterfly branch
                a[k + j] = u + v       // in-place update
                a[k + j + m/2] = u - v
                omega = omega * omega_m    // update to next twiddle factor
\end{verbatim}

And the inverse FFT can be implemented similarly, with the main difference being
the use of the inverse twiddle factors and a final scaling step to divide by \( n \).

\subsection{Inverse FFT (IFFT)}

The inverse FFT transforms the frequency-domain representation back to the time domain.
The mathematical definition of the inverse DFT is:
\[
a_k = \frac{1}{n} \sum_{j=0}^{n-1} A_j \cdot \omega_n^{-jk}
\]
where \(\omega_n^{-1} = e^{2\pi i / n}\) is the inverse primitive root of unity.

\subsubsection{Key Differences from FFT}

\begin{enumerate}
    \item \textbf{Inverse twiddle factors}: Use \(\omega_m^{-1} = e^{2\pi i / m}\) instead of \(\omega_m = e^{-2\pi i / m}\)
    \item \textbf{Final scaling}: Divide all results by \(n\) after the butterfly operations
\end{enumerate}

\subsubsection{Step-by-Step Example: IFFT for \(n=8\)}

Given the DFT result \([A_0, A_1, A_2, A_3, A_4, A_5, A_6, A_7]\), we first apply bit-reversal permutation:
\[
[A_0, A_4, A_2, A_6, A_1, A_5, A_3, A_7]
\]

\textbf{Stage 1: \(m=2\) (groups of size 2)}

With inverse twiddle factor \(\omega_2^0 = 1\):
\begin{align*}
\ &[A_0 + A_4, \; A_0 - A_4, \; A_2 + A_6, \; A_2 - A_6, \; A_1 + A_5, \; A_1 - A_5, \; A_3 + A_7, \; A_3 - A_7] \\
\ &= [B_0^{(1)}, B_4^{(1)}, B_2^{(1)}, B_6^{(1)}, B_1^{(1)}, B_5^{(1)}, B_3^{(1)}, B_7^{(1)}]
\end{align*}

\textbf{Stage 2: \(m=4\) (groups of size 4)}

Inverse twiddle factors: \(\omega_4^0 = 1\), \(\omega_4^{-1} = i\)

For the first half:
\begin{align*}
B_0^{(2)} &= B_0^{(1)} + \omega_4^0 \cdot B_2^{(1)} \\
B_2^{(2)} &= B_0^{(1)} - \omega_4^0 \cdot B_2^{(1)} \\
B_4^{(2)} &= B_4^{(1)} + \omega_4^{-1} \cdot B_6^{(1)} \\
B_6^{(2)} &= B_4^{(1)} - \omega_4^{-1} \cdot B_6^{(1)}
\end{align*}

For the second half:
\begin{align*}
B_1^{(2)} &= B_1^{(1)} + \omega_4^0 \cdot B_3^{(1)} \\
B_3^{(2)} &= B_1^{(1)} - \omega_4^0 \cdot B_3^{(1)} \\
B_5^{(2)} &= B_5^{(1)} + \omega_4^{-1} \cdot B_7^{(1)} \\
B_7^{(2)} &= B_5^{(1)} - \omega_4^{-1} \cdot B_7^{(1)}
\end{align*}

\textbf{Stage 3: \(m=8\) (groups of size 8, final stage)}

Inverse twiddle factors: \(\omega_8^0 = 1\), \(\omega_8^{-1} = \frac{\sqrt{2}}{2} + i\frac{\sqrt{2}}{2}\), \(\omega_8^{-2} = i\), \(\omega_8^{-3} = -\frac{\sqrt{2}}{2} + i\frac{\sqrt{2}}{2}\)

\begin{align*}
B_0^{(3)} &= B_0^{(2)} + \omega_8^0 \cdot B_1^{(2)} \\
B_1^{(3)} &= B_0^{(2)} - \omega_8^0 \cdot B_1^{(2)} \\
B_2^{(3)} &= B_2^{(2)} + \omega_8^{-1} \cdot B_3^{(2)} \\
B_3^{(3)} &= B_2^{(2)} - \omega_8^{-1} \cdot B_3^{(2)} \\
B_4^{(3)} &= B_4^{(2)} + \omega_8^{-2} \cdot B_5^{(2)} \\
B_5^{(3)} &= B_4^{(2)} - \omega_8^{-2} \cdot B_5^{(2)} \\
B_6^{(3)} &= B_6^{(2)} + \omega_8^{-3} \cdot B_7^{(2)} \\
B_7^{(3)} &= B_6^{(2)} - \omega_8^{-3} \cdot B_7^{(2)}
\end{align*}

\textbf{Final Scaling:} Divide all results by \(n = 8\):
\[
a_k = \frac{B_k^{(3)}}{8} \quad \text{for } k = 0, 1, \ldots, 7
\]

\subsubsection{General In-Place IFFT Pseudocode}

\begin{verbatim}
IFFT(a, n):
    // Step 1: Bit-reversal permutation (same as FFT)
    for i = 0 to n-1:
        j = bit_reverse(i, log2(n))
        if i < j:
            swap(a[i], a[j])

    // Step 2: Iterative butterfly operations with INVERSE twiddle factors
    for m = 2; m <= n; m *= 2:        // m = current group size
        omega_m = exp(2*pi*i / m)     // INVERSE principal twiddle factor
        for k = 0; k < n; k += m:     // iterate over groups
            omega = 1                  // omega_m^0
            for j = 0; j < m/2; j++:  // butterfly pairs within group
                u = a[k + j]           // upper butterfly branch
                v = omega * a[k + j + m/2]  // lower butterfly branch
                a[k + j] = u + v       // in-place update
                a[k + j + m/2] = u - v
                omega = omega * omega_m    // update to next twiddle factor

    // Step 3: Final scaling (divide by n)
    for i = 0 to n-1:
        a[i] = a[i] / n
\end{verbatim}

\subsubsection{Relationship between FFT and IFFT}

The IFFT can also be computed using the same FFT routine with a clever trick:
\begin{enumerate}
    \item Take the complex conjugate of the input: \(\overline{A_j}\)
    \item Apply the standard FFT
    \item Take the complex conjugate of the result
    \item Divide by \(n\)
\end{enumerate}

This works because:
\[
\text{IFFT}(A)_k = \frac{1}{n} \sum_{j=0}^{n-1} A_j \omega_n^{-jk} = \frac{1}{n} \overline{\sum_{j=0}^{n-1} \overline{A_j} \omega_n^{jk}} = \frac{1}{n} \overline{\text{FFT}(\overline{A})_k}
\]

This implementation is more elegant and efficient, as it avoids the overhead of recursion 
and uses only \( O(1) \) additional space for temporary variables during the butterfly operations. 
The overall time complexity remains \( O(n \log n) \).

\subsection{NTT} 

The Number Theoretic Transform (NTT) is a variant of the FFT that operates in a finite 
field, typically modulo a prime number. Similar to the FFT, the NTT use primitive roots 
of unity in the finite field to avoid floating-point operations. 

Accoding the Fermat little theorem, for a prime number $p$ and an integer $g$ not divisible by $p$, we have:
$$
g^{p-1} \equiv 1 \mod p
$$
This somehow provides us a way to substitute the complex roots of unity with integer 
roots of unity in the finite field. The main problem is to find a appropriate $p$ and 
a smaller $g$ to implement the NTT. $g$ should not be too large to avoid overflow, 
and $p$ should be large enough to accommodate the results of the convolution without 
modular reduction affecting the correctness.

In practice, we often choose a prime of the form \( p = k \cdot 2^m + 1 \) to ensure 
that there are enough roots of unity for the NTT. Usually, we choose $g = 3$ and 
$$p = 998244353 = 119 \cdot 2^{23} + 1$$
which is a commonly used prime in competitive programming. This ensure that we can 
perform NTT for sequences of length up to \( 2^{23} \) without overflow issues. By 
substituting the complex roots of unity with the integer roots of unity in the finite 
field, we can perform convolution operations efficiently using NTT, 
similar to how we use FFT for real-valued sequences.

\section{My favourite application of FFT}  

\subsection{Shor's Algorithm} 

In quantum computing, Shor's algorithm is a quantum algorithm for integer factorization,
 which can efficiently factor large integers. 
 The algorithm relies on the quantum Fourier transform (QFT), which is a quantum 
 analogue of the classical FFT. The QFT allows us to find the period of a function, 
 which is a crucial step in Shor's algorithm for finding the factors of an integer.

In the RSA cryptosystem, the security relies on the difficulty of factoring large 
integers. This problem can be reduced to finding the period of a function defined as:
$$
f(x) = a^x \mod N
$$
where \( N \) is the integer we want to factor, and \( a \) is a randomly chosen 
integer less than \( N \). If we can find a $r$ such that $f(x) = 1$ and $r$ is even, 
$$
(a^{r/2} - 1)(a^{r/2} + 1) \equiv 0 \mod N
$$
Then, if $a^{r/2 + 1} \not \equiv 0 \mod N$, then 
$$
\gcd(a^{r/2} - 1, N) \text{ and } \gcd(a^{r/2} + 1, N)
$$
are non-trivial factors of \( N \). The QFT is used to find the period \( r \) efficiently,
which is the key to breaking the RSA encryption. This application of FFT in quantum computing
is a groundbreaking development that has significant implications for cryptography and computational 
complexity.

The algorithm can be summarized as follows:
\begin{enumerate}
    \item First, we prepare an input $x$ to the register, which is a superposition of all possible values. 
    \item Choose a random $a$ for the function $f$. Then we apply the function $f(x)$ to the register, and put the result in another register. 
    \item We observe the second register and it collapses to a specific value $f(x_0)$. The first register is now in a
          superposition state which can be expressed as: $$\mid \psi \rangle = \sum_{k=0}^\infty \mid x_0 + kr \rangle$$ 
    \item We apply the QFT to the first register to transform $f = \sum_{x\in \mathbb{Z}_{2^n}} f(x) \mid x\rangle$ to 
          $\hat{f} = \sum_{u\in \mathbb{Z}_{n}} \hat{f}(x) \mid x\rangle$, where $\hat{f}(x) = \frac{1}{\sqrt{2^n}} \sum_{y \in \mathbb{Z}_{2^n}} \omega^{xy} f(x)$. 
    \item Then, we measure the first register, and it collapses to a specific value $u$. With high probability, $u$ will be close to a multiple of $\frac{2^n}{r} \cdot j, j\in \mathbb{Z}$, 
          which allows us to estimate the period $r$ using continued fractions.
\end{enumerate}

In step 4, the QFT is crucial for transforming the superposition state into a form that allows us to 
extract information about the period \( r \). We omit the details of the QFT implementation here, 
but it is very similar to the classical FFT that we have discussed in the previous sections. 
Here, we substitute the arithmetic operations with quantum gates, and the overall structure of the 
algorithm remains the same.

Shor's algorithm is a trade-off between accuracy and efficiency. We may find multiple $r$ that are not 
the actual period, but it can be proved that with high probability (a constant), 
we can find the correct period after a few iterations.

In digital signal processing, the FFT is used for analyzing the frequency components of signals. 
This is similar to how the QFT is used in Shor's algorithm to analyze the periodicity of a function. 

\subsection{Finite Impulse Response (FIR) Filter} 

In signal processing, FIR filters are used to process signals by convolving the input signal with a 
finite impulse response. The filter output a new signal that is a weighted sum of the current and previous input values. 
The convolution operation can be efficiently computed using the FFT, which allows us to perform the filtering in the frequency domain.

Formally, given an input signal \( x[n] \) and an FIR filter with impulse response \( h[n] \), the output signal \( y[n] \) is given by the convolution:
\[
y[n] = (x * h)[n] = \sum_{k=0}^{M-1} h[k] \cdot x[n-k]
\]
where \( M \) is the length of the impulse response, and $h[k]$ is the coefficient of the impulse response polynomial $H(z) = \sum_{k=0}^{M-1} h[k] z^{-k}$, $x[n-k]$ is the 
input signal delayed by $k$ samples in the sliding window $X(z) = \sum_{n=0}^{N-1} x[n] z^{-n}$ ($z^{-n}$ is called the Z-transform), but this is nothing different from the polynomial  
scenario in the previous FFT. To compute this convolution efficiently, we can use the FFT as follows:
\begin{enumerate}
    \item Compute the FFT of the input signal \( X[k] = \text{FFT}(x[n]) \) and the impulse response \( H[k] = \text{FFT}(h[n]) \).
    \item Perform point-wise multiplication in the frequency domain: \( Y[k] = X[k] \cdot H[k] \).
    \item Compute the inverse FFT to obtain the output signal: \( y[n] = \text{IFFT}(Y[k]) \).
\end{enumerate}

\subsection{Image Compression} 

The FFT is also widely used in image processing, particularly in image compression techniques such as JPEG.
In JPEG compression, the image is divided into 8x8 blocks of pixels, and the 2D FFT is applied to each block to transform the pixel values into the frequency domain.
The resulting frequency coefficients are then quantized, which reduces the precision of the coefficients and allows for more efficient storage. 
The quantization step is crucial for achieving high compression ratios, as it discards less important frequency components while retaining the most significant ones.

We formalize the 2D FFT as follows: Given an 8x8 block of pixel values \( P[i][j] \), we compute the 2D 
FFT to obtain the frequency coefficients \( F[u][v] \):
\[
F[u][v] = \sum_{i=0}^{7} \sum_{j=0}^{7} P[i][j] \cdot e^{-2\pi i \left( \frac{ui}{8} + \frac{vj}{8} \right)}
\]
where \( u, v = 0, 1, \ldots, 7 \). The resulting coefficients \( F[u][v] \) represent the frequency components of the image block.

After quantization, the coefficients are further compressed using techniques such as run-length encoding and Huffman coding to achieve even higher compression ratios.
In practice, the JPEG compression use Discrete Cosine Transform (DCT) instead of the FFT, but the underlying 
principle of transforming the image into the frequency domain and quantizing the coefficients is similar to the FFT-based approach. 
The only difference is that the DCT uses cosine functions as convolution kernels. 

\subsection{FlashFFTConv \cite{fu2023flashfftconvefficientconvolutionslong}}

\subsection{High Definition Radar \cite{rebut2022rawhighdefinitionradarmultitask}}

\section{Conclusion}

\end{document}
